\documentclass{article}
\usepackage[T1]{fontenc}
\usepackage{lmodern}
\usepackage{graphicx}
\usepackage{amsmath,amssymb}
\usepackage[a4paper,margin=2cm]{geometry}
\usepackage[absolute,overlay]{textpos}
\setlength{\TPHorizModule}{1mm}
\setlength{\TPVertModule}{1mm}
\usepackage[indonesian]{babel}
\usepackage{hyperref}
\setlength{\emergencystretch}{2em}
% unit: Halaman 42

\begin{document}

% === Halaman 42 (python/native) ===

Algoritma Elliptic Curve Diffie-Hellman  (ECDH)

• Alice dan Bob ingin berbagi sebuah kunci rahasia K yang sama.

• Alice dan Bob menyepakati kurva eliptik y2 = x3 + ax + b  mod p dan titik B pada kurva

• Alice dan Bob menghitung kunci publik dan kunci privat masing-masing. 

• Alice

• Kunci privat = a

• Kunci publik = PA = a  B

• Bob

• Kunci privat = b

• Kunci publik = PB = b  B

• Alice dan Bob saling mengirim kunci publik masing-masing.

• Keduanya melakukan perkalian kunci privatnya dengan kunci publik mitranya untuk 

mendapatkan kunci rahasia yang mereka bagi

• Alice → KAB = aPB = a(b  B)

• Bob  → KAB = bPA = b(a  B)

• Kunci rahasia = KAB = abB

*) Sumber bahan: Debdeep Mukhopadhyay, Elliptic Curve Cryptography ,

 Dept of Computer Sc and Engg IIT Madras

42

\end{document}
